% === Begin Label Data ===
% ID: 276
% 难度星级: 
% 标签: 
% 备注: 
% 组卷引用次数: 0
% === End  Label Data ===

\begin{problem}{2026}{M}{深圳中学2026届高三二轮二阶测试}{19}{概率}
数学中的等效性与对称性不仅体现在几何和函数中，还体现在概率问题中．已知一个不透明的罐子中有编号分别为 $1,2,3,\cdots,2n$ 的 $2n$ 个小球（$n\geq2$），除编号外，小球的大小、质地完全相同．现从中一次性随机摸出 $k$ 个小球．

（1）若 $n=4$，$k=3$，求这 $k$ 个小球编号两两不相邻的概率；

（2）记 $X,Y$ 分别为摸出的这 $k$ 个小球的最大、最小编号，求 $E(X+Y)$；

（3）若 $k=n$，记 $\Delta$ 表示摸出小球的最大最小编号之差，$\bar{\Delta}$ 表示未摸出小球的最大最小编号之差，记 $P_n(A)=P(\Delta=\bar{\Delta})$，证明 $P_n(A)>\dfrac{1}{6}\left(1-\dfrac{1}{4^{n-2}}\right)$.

参考公式：已知随机变量的期望具有线性可加性，即对于随机变量 $\xi,\eta$，有 $E(\xi\pm\eta)=E(\xi)\pm E(\eta)$.
\end{problem}